\section{From small defect to a common support region}\label{short:acq:section}

The localization theorem bounds $nD$ and $|H|$.
Lemma~\ref{ms:joint-defect}
bounds $D$ below by nonnegative costs, so a small defect makes each
cost small: the sign probability must be close to $p_*$, the support
gap must nearly reach the conditional means, and the spread within each
sign class must be small. Proposition~\ref{short:acq:geometry} makes
these conclusions quantitative using the defect alone.
Proposition~\ref{short:acq:strip} then uses contact and the sharp
predecessor bound to control the entire conditional support. Together
with the threshold bound, these conclusions give the domain needed by
Theorem~\ref{short:count:local}.
Throughout this section set
\[
 \sigma=\sqrt{10},\qquad p_*=(4-\sigma)/2,\qquad
 B_*=32.20665378.
\]
The law, support gaps and sign split are those of
Section~\ref{ms:geometry-counts}; in particular $D$ is the actual gap
defect, not a moment surrogate. Write the central gap, once acquired, as
$(-A_-,A_+)$, using the notation of
Section~\ref{ms:geometry-counts}. The first use of $n\ge700$ is the
inequality $B_*/700<.047$ below; the count comparison also uses this
sample-size bound in its error estimates.

\begin{proposition}[Direct geometry]\label{short:acq:geometry}
Let $X$ be compactly supported, centered and of variance one. If
$n\ge700$ and $nD\le B_*$, its closed support has a gap
$(-A_-,A_+)$ across zero, with $A_-,A_+>0$. In the actual decomposition
$X=\ell(B_0+U-p)$, put
\[
 \begin{gathered}
a_0=p-A_-/\ell,\quad a_1=q-A_+/\ell,\quad d=a_0+a_1,\\
\theta=\E\{X-\E(X\mid B_0)\}^2,\qquad
\lambda=n\E U^2=n\theta/\ell^2.
\end{gathered}
\]
Then $a_0,a_1\ge0$, $\E(U\mid B_0)=0$, $U_0\le a_0$,
$U_1\ge-a_1$, and
\begin{gather}
 \frac38<p<\frac{93}{200},\qquad
 \theta<.0036,\qquad d<.004,\qquad\ell>1.996,
 \label{short:acq:domain}\\
 \rho<\frac98,\qquad .066148<\beta<.2527,
 \qquad\lambda<\frac7{20}.\label{short:acq:outputs}
\end{gather}
These conclusions require no contact or predecessor hypothesis.
\end{proposition}

The estimates are obtained in the following order. Each preliminary
bound supplies a hypothesis for a sharper one later in the proof.
\begin{center}\small
\begin{tabular}{@{}p{.27\linewidth}p{.65\linewidth}@{}}
\toprule
Input & Purpose of the output \\
\midrule
$D<.047$ and the zero-in-support bound & Establish a central gap and
nonnegative inward displacements. \\
The joint defect budget & Obtain $.36<p<.48$, $\theta<.006$ and
$\ell>1.99$ before taking derivatives of the refined bounds. \\
The inward-motion term & Give $d<.004$. \\
Monotonicity in $d$ and $\theta$ & Give the final probability interval,
$\theta<.0036$ and $\ell>1.996$. \\
The conditional cubic costs & Control the total count variance
$\lambda$, rather than only the variance of $X$. \\
Contact and the predecessor bound & Exclude the two outer support rays
in Proposition~\ref{short:acq:strip}. \\
\bottomrule
\end{tabular}
\end{center}

\begin{proof}
\emph{1. A central gap and preliminary bounds.}
We have $D<D_0:=.047$, because $700D_0>B_*$. If zero were in the
closed support, Lemma~\ref{ms:physical-defect} would give
$D\ge(\sigma-1)\rho\ge\sigma-1>D_0$. Thus the central gap exists;
centering supplies both signs. Conditional means imply $a_j\ge0$ and
$0\le d<1$. Put $v=pq$ and $w=1-\theta=v\ell^2$, and use the
probability polynomial $\mathcal Q$ from
Section~\ref{ms:geometry-counts}.
We first retain all the costs of the actual conditional noises. Let
\[
 \int f\,d\mu_\pm
 =\E\bigl[(X-\E(X\mid B_0))^2\ind_{\{\pm X>0\}}f(|X|)\bigr].
\]
Their masses sum to $\theta$. In this notation the joint budget
\eqref{ms:joint-defect-budget} reads
\begin{equation}\label{short:acq:budget}
 D\ge\ell w\mathcal Q+E_0+\int k_+\,d\mu_++\int k_-\,d\mu_-,
\end{equation}
where
\begin{gather*}
 E_0=6\ell w\{d(1-d)+a_0a_1\},\\
 k_+(x)=(\sigma-1)(x+2\ell q)+6A_+,\qquad
 k_-(x)=(\sigma+1)(x+2\ell p)+6A_-.
\end{gather*}
The measures $\mu_\pm$ abbreviate the displayed variance-weighted
expectations under the same marginal $P$.

Hereafter use $3.162277660<\sigma<3.162277661$. The weaker gap estimate
$J\le\rho+2mA_-A_+$, together with $\rho\ge1$, $\gamma\le\rho$ and
$m\le1/2$, gives
\[
 A_-A_+\ge z:=\frac{3.162277660-1-D_0}{3}>0,
 \qquad w\ge A_-A_+\ge z.
\]
The second inequality follows from $A_-\le\ell p$, $A_+\le\ell q$.
All costs in \eqref{short:acq:budget} are nonnegative, so
\[
 D\ge\frac{w^{3/2}\mathcal Q}{\sqrt v}\ge2z^{3/2}\mathcal Q.
\]
Let $\mathcal Q_-$ replace $\sigma$ in its polynomial expression by the lower
rational bound. The two exact tests
\[
 4z^3\mathcal Q_-(.36)^2>D_0^2,\qquad
 4z^3\mathcal Q_-(.48)^2>D_0^2
\]
exclude $p\le.36$ and $p\ge.48$, respectively: the squared distance
to $p_*\in(.36,.48)$ is monotone outside that interval.

\emph{2. Preliminary noise and inward-displacement bounds.}
On $.36<p<.48$, $\ell\ge2\sqrt z>1.67$, and the noise coefficients satisfy
\begin{align*}
 k_+(x)/\ell&\ge q\{(\sigma+5)z/w+2(\sigma-1)\}>5.2,\\
 k_-(x)/\ell&\ge p\{(\sigma+7)z/w+2(\sigma+1)\}>5.2.
\end{align*}
For the first use $x\ge A_+$ and $A_+\ge z/(\ell p)$; for the second
use their negative counterparts. Since $w\le1$, the constants follow
by taking $q=.52$ and $p=.36$. Therefore
$D\ge1.67(5.2)\theta$, whence $\theta<.006$, $w>.994$ and
$\ell>1.99$. Also $A_-+A_+\ge2\sqrt z$ gives
\[
 d\le1-2\sqrt{zv/w}<.2,\qquad v\ge.36(.64).
\]
Evaluating the kernels at their lower support endpoints in
\eqref{short:acq:budget} now yields
\begin{equation}\label{short:acq:direct}
 D\ge\ell\{w\mathcal Q+6wd(1-d)
            +\theta\min(C_+-b_+d,C_--b_-d)\},
\end{equation}
where
\[
 C_+=3(\sigma+1)q,\quad C_-=3(\sigma+3)p,\quad
 b_+=\sigma+5,\quad b_-=\sigma+7.
\]
Each branch increases in $d$ on the current domain, since
$6(.994)(.6)-.006(3.162277661+7)>0$; both noise prices are positive.
The inward cost alone excludes $d\ge.004$, because
\[
 1.99(.994)6(.004)(.996)>.047.
\]

\emph{3. The final probability and variance bounds.}
After replacing $d$ by zero, the derivative in $\theta$ of either branch
$\sqrt{(1-\theta)/v}\{(1-\theta)\mathcal Q+\theta C\}$ has the sign of
$(2-3\theta)C-3(1-\theta)\mathcal Q$. It is positive: $C>6.49$, $\mathcal Q<.05$
and $\theta<.006$. Hence $D\ge\mathcal Q/\sqrt v$.
Since $\sqrt v\le1/2$, it follows that $D\ge2\mathcal Q(p)$.
Moreover
\[
 \frac38<p_*<\frac{21}{50}
      =\frac12\left(\frac38+\frac{93}{200}\right).
\]
Every $p$ outside $(3/8,93/200)$ is therefore at least as far from
$p_*$ as $3/8$, so $\mathcal Q(p)\ge \mathcal Q(3/8)$. The single rational comparison
\[
 \mathcal Q_-(3/8)=\frac{37936011}{1600000000}
       >\frac{47}{2000}=\frac{D_0}{2}
\]
contradicts $D<D_0$ on both outside rays, including the endpoints.
Thus $3/8<p<93/200$.
Since
\[
 C\ge\min\{3(\sigma+1)(107/200),\;3(\sigma+3)(3/8)\},
\]
the bound $D\ge\ell\theta C$, with $\ell>1.99$, improves
$\theta<.0036$ and consequently $\ell>1.996$.
Here the $d$-dependent terms may be dropped together, since
$6w(1-d)-\theta b_\pm>0$ on the current domain. This keeps the displacement cost together with its signed noise correction.

\emph{4. The remaining moment bounds.} The exact conditional cubic expansion gives
\[
 \rho=w^{3/2}\rho_B(p)+\int(x+2\ell q)\,d\mu_+
                          +\int(x+2\ell p)\,d\mu_-,
 \qquad\rho_B(p)=\frac{1-2v}{\sqrt v}.
\]
Because the two integrands are priced by at least $\sigma-1$ in
\eqref{short:acq:budget},
\[
 \rho\le\rho_B(p)+D/(\sigma-1)
       <1.1+.047/(54/25)<9/8.
\]
Here $v\ge15/64$ implies $\rho_B<1.1$ by squaring positive quantities.
The signed-second-moment identity gives
\[
 1-2p-2q\theta\le\beta\le1-2p+2p\theta;
\]
its endpoints imply $.066148<\beta<.2527$.

\emph{5. From conditional spread to total count variance.}
A bound on $\theta$ alone would not control $\lambda=n\theta/\ell^2$
uniformly in $n$. We therefore retain more of the cubic cost. A
one-sided moment inequality supplies a variance-squared term divided by
the inward displacement; pairing it with the displacement cost by
AM--GM will give a larger linear coefficient of $\lambda$.
If $Z$ is centered, $Z\le a$, $a>0$, and $t=\E Z^2$, then
\[
 0\le\E[(a-Z)(Z+t/a)^2]=at-\E Z^3-t^2/a.
\]
Apply this to $U_0$ and $-U_1$. A zero endpoint forces the corresponding
noise to vanish, so no division by zero is needed. For
$\lambda_0=nq\E U_0^2$, $\lambda_1=np\E U_1^2$,
the same gap-budget expansion, with these two cubic bounds inserted, gives
\begin{align*}
 nD/\ell^3\ge{}&nv\mathcal Q+6nv\{d(1-d)+a_0a_1\}\\
 &+\{3(\sigma+3)p-(\sigma+7)a_0\}\lambda_0
  +\{3(\sigma+1)q-(\sigma+5)a_1\}\lambda_1\\
 &+\frac{(\sigma+1)\lambda_0^2}{nqa_0}
  +\frac{(\sigma-1)\lambda_1^2}{npa_1}.
\end{align*}
Terms with $a_j=0$ are absent. Set $\bar d=.004$ and use
$d(1-d)+a_0a_1\ge(1-\bar d)(a_0+a_1)$.
The reciprocal term is essential: pairing its variance-squared cost
with the inward-displacement cost gives a linear cost in $\lambda$.
Specifically, AM--GM gives
\[
 6nv(1-\bar d)a_0+
 \frac{(\sigma+1)\lambda_0^2}{nqa_0}
 \ge2\lambda_0\sqrt{6(1-\bar d)(\sigma+1)p}.
\]
Thus
\begin{equation}\label{short:acq:cubic-budget}
 nD\ge\ell^3\{nv\mathcal Q+\lambda_0Y_0(p)+\lambda_1Y_1(p)\},
\end{equation}
where
\begin{align*}
 Y_0&=3(\sigma+3)p-(\sigma+7)\bar d
                   +2\sqrt{6(1-\bar d)(\sigma+1)p},\\
 Y_1&=3(\sigma+1)q-(\sigma+5)\bar d
                   +2\sqrt{6(1-\bar d)(\sigma-1)q}.
\end{align*}
Each increases in its respective argument $p$ or $q$. At $p=3/8$ and
$q=107/200$, lower bounds $3$ and $2.6$ for the roots give
$Y_0>12.8$, $Y_1>11.8$. Since $\lambda_0+\lambda_1=\lambda$,
\[
 B_*\ge nD>1.996^3(11.8)\lambda
 \quad(\lambda>0),\qquad
 1.996^3(11.8)(.35)>B_*.
\]
This proves $\lambda<.35$; at zero noise it is immediate.
\end{proof}

\begin{proposition}[Two contact rays]\label{short:acq:strip}
Under the hypotheses of Proposition~\ref{short:acq:geometry}, suppose
in addition that $C_{n-1}\le\CE$ and that $(P,H)$ maximizes the excess
at $\CE$, with nonnegative maximum. Then
\[
                -.6<U_0\le a_0,\qquad -a_1\le U_1<.6.
\]
The support conclusion is uniform in the raw threshold $H$.
\end{proposition}

\begin{proof}
A support point too far from its conditional mean would make the
contact upper expression \eqref{v9:entry:contact} negative. The
all-threshold Gaussian heat estimate and its radial derivative bound
from the proof of Theorem~\ref{v9:entry:initial} give, for $t>\sqrt3$,
\begin{align*}
 H_n(H,\pm t)&\le \overline H(t):=S^3\left\{\frac25\left(
 \frac{t^3-3t}{6}+\frac{|t^4-6t^2+3|}{72}\right)
 +\frac{t^5}{4100}+\frac{3t}{820}
 +\frac{17(3t^2-2)}{78720}\right\},\\
 |\partial_t H_n(H,\pm t)|&<g_{\rm rad}(t):=\frac25(t^2/2+1/2+\eta),
 \qquad S=81/80,\quad\eta=8/875.
\end{align*}
The first bound follows directly from \eqref{v9:entry:heat}:
$\|\phi'''\|_\infty<2\phi(0)$, $n-1\ge41$, and
$\sqrt{n-1}>6$. No derivative of the absolute-value majorant is taken.

At a support point of magnitude $t$, contact requires a nonnegative upper
expression $H_n(H,\pm t)+\CE Q_b(t)$, where
\[
 Q_b(t)=-t^3+\tfrac32rt^2+3bt+\eta r,
 \quad r=9/8,
 \quad b=\begin{cases}-.066148&\text{negative ray},\\
                         .2527&\text{positive ray}.
          \end{cases}
\]
These replacements follow with their signs from
\eqref{short:acq:outputs}; the discarded excess term is nonpositive.
Take $c_0=819/2000<\CE$ and the two fixed boundaries
\[
 t_-=1.996(3/8+.6)=1.9461,
 \qquad t_+=1.996(107/200+.6)=2.26546.
\]
For a fixed raw threshold $H$, abbreviate the appropriate radial
function $H_n(H,\pm t)$ by $H_n(t)$, with derivative taken in the
positive magnitude $t$. Set
\[
 \Xi(t)=g_{\rm rad}(t)+c_0Q_b'(t).
\]
Exact rational substitution gives the following conservative bounds:
\begin{center}
\begin{tabular}{@{}lcc@{}}
\toprule
Boundary quantity & $t_-$, negative ray & $t_+$, positive ray \\
\midrule
$\overline H(t)+c_0Q_b(t)$ & $<-2/5$ & $<-13/100$ \\
$\Xi(t)$ & $<-1$ & $<-8/5$ \\
$\Xi'(t)$ & $<-13/5$ & $<-16/5$ \\
\bottomrule
\end{tabular}
\end{center}
Both boundaries satisfy $t^2>3$. The constant second derivative is
\[
 \Xi''=\frac25-6c_0=-\frac{2057}{1000}<0.
\]
Consequently $\Xi'$ and then $\Xi$ remain negative on each outward ray.
Because $g_{\rm rad}>0$, the identity $Q_b'=(\Xi-g_{\rm rad})/c_0$ gives $Q_b'<0$ there.
At the boundary $\overline H\ge0$, so the first table entry gives $Q_b<0$;
this sign also persists outward. Since $\CE>c_0$, we obtain
\[
 H_n(t_\pm)+\CE Q_b(t_\pm)<0,\qquad
 H_n'(t)+\CE Q_b'(t)\le g_{\rm rad}(t)+c_0Q_b'(t)=\Xi(t)<0.
\]
The contact upper expression therefore cannot be nonnegative anywhere
on either ray. This excludes every support point there, including
boundary atoms. We differentiate $H_n$ itself; the absolute-value
majorant $\overline H$ is used only at the two boundary points.

On the negative class, $X>-t_-$ and $\ell>1.996$ give
$U_0=X/\ell+p>-(3/8+.6)+p>-.6$. On the positive class,
$X<t_+$ gives
$U_1=X/\ell+p-1<107/200+.6+p-1<.6$.
The inward endpoints and conditional centering are the actual ones from
Proposition~\ref{short:acq:geometry}.
\end{proof}

By \eqref{ms:bernoulli-coordinates}, the law compared in the next
section is exactly the original law $X$, with the same threshold $H$.
The sign probability and the signed moments refer throughout to this
chosen law.
